% GENERATED FILE. Edit canonical lessons, then run npm run generate:books.
% canonical-source-hash: fnv1a64:e1a33d8e7637b14f
% canonical-lessons: TE-C242-teccu, TE-C242-upayogincu, TE-C242-mugincu, TE-C242-modalupettu, TE-C242-istri-ceyu

\chapter{To bring, To use, To finish, To begin, To iron}
\label{ch:te-to-bring-to-use-to-finish-to-begin-to-iron}
\hlchaptermodality{242}

\begin{chapteropening}
\textbf{By the end of this chapter:} \emph{I can say to bring, to use, to finish, to begin and to iron.}

Complete the last lesson of chapter 242: I can say to bring, to use, to finish, to begin and to iron.
\end{chapteropening}

\section[\texorpdfstring{\te{తెచ్చు} (teccu)}{తెచ్చు (teccu)}]{\te{తెచ్చు} (teccu) — to bring}

\label{lesson:TE-C242-teccu}



\begin{quote}
Before the new one: say the Telugu for to sneeze, then the Telugu for to vomit.
\end{quote}

\subsection*{You'll want to know: \te{తెచ్చు}}
\textbf{\te{తెచ్చు}} — \emph{teccu} — \textquotedblleft{}to bring\textquotedblright{}.

\begin{grammarlens}[title={What you've built}]
One more word for this chapter.
\end{grammarlens}

\subsection*{Your turn}
\begin{itemize}
\raggedright
  \item \emph{Say it:} \emph{teccu}
  \item \emph{Say it:} \emph{teccu}, once more
  \item \emph{Say it:} say \emph{teccu}
  \item \emph{Recall:} say the Telugu for to sneeze, then the Telugu for to vomit, then say \emph{teccu} again
\end{itemize}

\begin{tcolorbox}[breakable,colback=teal!4,colframe=teal!35!black,title={Before you move on}]
What does \te{తెచ్చు} mean? (\textquotedblleft{}To bring\textquotedblright{}.) Say it once more.
\end{tcolorbox}
\section[\texorpdfstring{\te{ఉపయోగించు} (upayōgiñcu)}{ఉపయోగించు (upayōgiñcu)}]{\te{ఉపయోగించు} (upayōgiñcu) — to use}

\label{lesson:TE-C242-upayogincu}



\begin{quote}
Before the new one: say the Telugu for to vomit, then the Telugu for to bring.
\end{quote}

\subsection*{You'll want to know: \te{ఉపయోగించు}}
\textbf{\te{ఉపయోగించు}} — \emph{upayōgiñcu} — \textquotedblleft{}to use\textquotedblright{}.

\begin{grammarlens}[title={What you've built}]
One more word for this chapter.
\end{grammarlens}

\subsection*{Your turn}
\begin{itemize}
\raggedright
  \item \emph{Say it:} \emph{upayōgiñcu}
  \item \emph{Say it:} \emph{upayōgiñcu}, once more
  \item \emph{Say it:} say \emph{upayōgiñcu}
  \item \emph{Recall:} say the Telugu for to vomit, then the Telugu for to bring, then say \emph{upayōgiñcu} again
\end{itemize}

\begin{tcolorbox}[breakable,colback=teal!4,colframe=teal!35!black,title={Before you move on}]
What does \te{ఉపయోగించు} mean? (\textquotedblleft{}To use\textquotedblright{}.) Say it once more.
\end{tcolorbox}
\section[\texorpdfstring{\te{ముగించు} (mugiñcu)}{ముగించు (mugiñcu)}]{\te{ముగించు} (mugiñcu) — to finish}

\label{lesson:TE-C242-mugincu}



\begin{quote}
Before the new one: say the Telugu for to bring, then the Telugu for to use.
\end{quote}

\subsection*{You'll want to know: \te{ముగించు}}
\textbf{\te{ముగించు}} — \emph{mugiñcu} — \textquotedblleft{}to finish\textquotedblright{}.

\begin{grammarlens}[title={What you've built}]
One more word for this chapter.
\end{grammarlens}

\subsection*{Your turn}
\begin{itemize}
\raggedright
  \item \emph{Say it:} \emph{mugiñcu}
  \item \emph{Say it:} \emph{mugiñcu}, once more
  \item \emph{Say it:} say \emph{mugiñcu}
  \item \emph{Recall:} say the Telugu for to bring, then the Telugu for to use, then say \emph{mugiñcu} again
\end{itemize}

\begin{tcolorbox}[breakable,colback=teal!4,colframe=teal!35!black,title={Before you move on}]
What does \te{ముగించు} mean? (\textquotedblleft{}To finish\textquotedblright{}.) Say it once more.
\end{tcolorbox}
\section[\texorpdfstring{\te{మొదలుపెట్టు} (modalupeṭṭu)}{మొదలుపెట్టు (modalupeṭṭu)}]{\te{మొదలుపెట్టు} (modalupeṭṭu) — to begin}

\label{lesson:TE-C242-modalupettu}



\begin{quote}
Before the new one: say the Telugu for to use, then the Telugu for to finish.
\end{quote}

\subsection*{You'll want to know: \te{మొదలుపెట్టు}}
\textbf{\te{మొదలుపెట్టు}} — \emph{modalupeṭṭu} — \textquotedblleft{}to begin\textquotedblright{}.

\begin{grammarlens}[title={What you've built}]
One more word for this chapter.
\end{grammarlens}

\subsection*{Your turn}
\begin{itemize}
\raggedright
  \item \emph{Say it:} \emph{modalupeṭṭu}
  \item \emph{Say it:} \emph{modalupeṭṭu}, once more
  \item \emph{Say it:} say \emph{modalupeṭṭu}
  \item \emph{Recall:} say the Telugu for to use, then the Telugu for to finish, then say \emph{modalupeṭṭu} again
\end{itemize}

\begin{tcolorbox}[breakable,colback=teal!4,colframe=teal!35!black,title={Before you move on}]
What does \te{మొదలుపెట్టు} mean? (\textquotedblleft{}To begin\textquotedblright{}.) Say it once more.
\end{tcolorbox}
\section[\texorpdfstring{\te{ఇస్త్రీ} \te{చేయు} (istrī cēyu)}{ఇస్త్రీ చేయు (istrī cēyu)}]{\te{ఇస్త్రీ} \te{చేయు} (istrī cēyu) — to iron (clothes)}

\label{lesson:TE-C242-istri-ceyu}



\begin{quote}
Before the new one: say the Telugu for to finish, then the Telugu for to begin.
\end{quote}

\subsection*{You'll want to know: \te{ఇస్త్రీ} \te{చేయు}}
\textbf{\te{ఇస్త్రీ} \te{చేయు}} — \emph{istrī cēyu} — \textquotedblleft{}to iron (clothes)\textquotedblright{}.

\begin{grammarlens}[title={What you've built}]
One more word for this chapter.
\end{grammarlens}

\subsection*{Your turn}
\begin{itemize}
\raggedright
  \item \emph{Say it:} \emph{istrī cēyu}
  \item \emph{Say it:} \emph{istrī cēyu}, once more
  \item \emph{Say it:} say \emph{istrī cēyu}
  \item \emph{Recall:} say the Telugu for to finish, then the Telugu for to begin, then say \emph{istrī cēyu} again
\end{itemize}

\begin{tcolorbox}[breakable,colback=teal!4,colframe=teal!35!black,title={Before you move on}]
What does \te{ఇస్త్రీ} \te{చేయు} mean? (\textquotedblleft{}To iron (clothes)\textquotedblright{}.) Say it once more.
\end{tcolorbox}
